\section{Bölüm 4: Belirli İntegral}

Bu dosya LaTeX etiketleyicisinin (\texttt{aicls.book.ingest.latex\_tagger})
tanıdığı ortamları gösterir. Her ortam bir \verb|\label| taşımak zorundadır.

\begin{tanim}[Riemann Toplamı]\label{tanim:4.1}
$[a, b]$ aralığının bir bölüntüsü için her alt aralıktan seçilen örnek noktalarla
kurulan toplama Riemann toplamı denir.
\[
S_n = \sum_{i=1}^{n} f(x_i^*) \, \Delta x_i
\]
\end{tanim}

\begin{tanim}[Belirli İntegral]\label{tanim:4.2}
Bölüntü normu sıfıra giderken Riemann toplamlarının limiti varsa bu limite
fonksiyonun belirli integrali denir.
\[
\int_a^b f(x)\,dx = \lim_{\|P\| \to 0} S_n
\]
\end{tanim}

\begin{teorem}[Analizin Temel Teoremi]\label{teorem:4.3}
$f$ fonksiyonu $[a, b]$ aralığında sürekli ve $F$ onun bir ilkeli ise belirli
integral uç noktalardaki ilkel değerlerinin farkına eşittir.
\[
\int_a^b f(x)\,dx = F(b) - F(a)
\]
\end{teorem}

\begin{ornek}[Alan Hesabı]\label{ornek:4.4}
Parabolün altında kalan alan belirli integralle hesaplanır.
\[
\int_0^1 x^2 \, dx = \frac{1}{3}
\]
\end{ornek}
